\documentclass{article}
\input{../lib.sty}

\title{\texttt{fn sum\_error}}
\author{\MichaelShoemateAuthor}

\begin{document}
\maketitle
\contrib

\section{Hoare Triple}
\subsection*{Precondition}
\subsubsection*{Compiler-verified}
\begin{itemize}
    \item \texttt{T: Float}, either \texttt{f32} or \texttt{f64}.
    \item $h=\texttt{max\_depth}$ is of type \texttt{usize}.
    \item $A=\texttt{max\_abs\_sum}$ and $B=\texttt{max\_norm}$ are of type \texttt{\&RBig}.
\end{itemize}

\subsubsection*{Caller-verified}
None.

\subsection*{Postcondition}
\begin{theorem}
    \label{postcondition}
    The function returns \texttt{Err(e)} or \texttt{Ok(out)}.
    In the latter case, for every $p\ge1$ and vector of sums
    \[
        s_j=\sum_i x_{ji},\qquad a_j=\sum_i |x_{ji}|,
        \qquad \|a\|_\infty\le A,\quad \|a\|_p\le B,
    \]
    evaluated by binary summation trees of depth at most $h$ with finite
    leaves of type $T$ and round-to-nearest addition with gradual underflow,
    all intermediate sums are finite and
    $\|\widehat{s}-s\|_p\le\texttt{out}$,
    where $\widehat{s}$ is the vector of computed sums.
\end{theorem}

\begin{lemma}\label{accuracy-bound}
    Let $u=\texttt{T::epsilon()}/2$, $hu<1$, and $\gamma_h=hu/(1-hu)$.
    For a nonoverflowing summation tree of depth at most $h$,
    \[
        \left|\widehat{s}-\sum_i x_i\right|\le\gamma_h\sum_i|x_i|.
    \]
\end{lemma}

\begin{proof}
    Each addition has $\operatorname{fl}(a+b)=(a+b)(1+\delta)$,
    $|\delta|\le u$ \href{https://doi.org/10.1137/0914050}{(Higham, 1993)}.
    Subnormal additions are exact.
    For the path $P_i$ from $x_i$ to the root, set
    $\theta_i=\prod_{k\in P_i}(1+\delta_k)-1$. Since $|P_i|\le h$,
    \begin{align*}
        |\theta_i|
        &\le \max\{(1+u)^h-1,1-(1-u)^h\}\\
        &\le (1-hu)^{-1}-1=\gamma_h,\\
        \left|\widehat{s}-\sum_i x_i\right|
        &=\left|\sum_i x_i\theta_i\right|
        \le\gamma_h\sum_i|x_i|.
    \end{align*}
\end{proof}

\begin{proof}[Proof of Theorem~\ref{postcondition}]
    Let $C=\min(A,B)\ge\|a\|_\infty$.
    If $h=0$ or $C=0$, the sum is exact and finite, and \texttt{out}=0.
    Otherwise, the exact rational checks ensure
    \[
        hu<1,\qquad C(1+\gamma_h)\le T_{\max},
    \]
    where $T_{\max}$ is the largest finite value of type $T$.

    At a hypothetical first overflowing addition, its child sums
    $\widehat{s}_L,\widehat{s}_R$ are finite.
    By Lemma~\ref{accuracy-bound}, with $a_L,a_R$ their absolute-summand sums,
    \begin{align*}
        |\widehat{s}_L+\widehat{s}_R|
        &\le(1+\gamma_h)(a_L+a_R)\\
        &\le(1+\gamma_h)C\le T_{\max},
    \end{align*}
    a contradiction. Thus the lemma applies to each coordinate:
    \[
        |\widehat{s}_j-s_j|\le\gamma_h a_j
        \implies \|\widehat{s}-s\|_p
        \le\gamma_h\|a\|_p\le\gamma_h B=\texttt{out}.
    \]
\end{proof}
\end{document}
